\documentclass[11pt]{article}
\usepackage[a4paper,margin=21mm]{geometry}
\usepackage[no-math]{fontspec}
\setmainfont{Latin Modern Roman}
\newfontfamily\CJKfont{Noto Serif CJK SC}
\XeTeXlinebreaklocale "zh"
\XeTeXlinebreakskip=0pt plus 1pt
\usepackage{amsmath,amssymb,amsthm,amscd,graphicx,color}
\usepackage{xurl}
\usepackage[unicode,hidelinks]{hyperref}
\allowdisplaybreaks
\setlength\emergencystretch{3em}
\numberwithin{equation}{section}

\usepackage[mathscr]{euscript}

\numberwithin{equation}{section}
\newtheorem{Theorem}{Theorem}[section]
\newtheorem{Corollary}[Theorem]{Corollary}
\newtheorem{Lemma}[Theorem]{Lemma}
\newtheorem{Proposition}[Theorem]{Proposition}
 { \theoremstyle{definition}
\newtheorem{Remark}[Theorem]{Remark} }

\newcommand{\supp}{\mathop {\mathrm {supp}}\nolimits}
\newcommand{\rk}{\mathop {\mathrm {rk}}\nolimits}
\newcommand{\Aut}{\mathop {\mathrm {Aut}}\nolimits}
\newcommand{\Out}{\mathop {\mathrm {Out}}\nolimits}
\renewcommand{\Re}{\mathop {\mathrm {Re}}\nolimits}
\newcommand{\res}{\mathop {\mathrm {res}}\nolimits}

\newcommand{\hyp}[2]{\fourIdx{}{#1}{}{\!#2}F}
\newcommand{\hypc}[2]{\fourIdx{}{#1}{\C}{\!#2}F}
\newcommand{\Kc}[2]{\fourIdx{}{#1}{\C}{\!#2}{\EuScript K}}
\newcommand{\Gc}[2]{\fourIdx{}{#1}{\C}{\!#2}G}

\def\ov{\overline}

\def\wt{\widetilde}

\renewcommand{\rk}{\mathop {\mathrm {rk}}\nolimits}
\renewcommand{\Aut}{\mathop {\mathrm {Aut}}\nolimits}
\renewcommand{\Re}{\mathop {\mathrm {Re}}\nolimits}
\renewcommand{\Im}{\mathop {\mathrm {Im}}\nolimits}
\newcommand{\sgn}{\mathop {\mathrm {sgn}}\nolimits}

\def\bfa{\mathbf a}
\def\bfb{\mathbf b}
\def\bfc{\mathbf c}
\def\bfd{\mathbf d}
\def\bfe{\mathbf e}
\def\bff{\mathbf f}
\def\bfg{\mathbf g}
\def\bfh{\mathbf h}
\def\bfi{\mathbf i}
\def\bfj{\mathbf j}
\def\bfk{\mathbf k}
\def\bfl{\mathbf l}
\def\bfm{\mathbf m}
\def\bfn{\mathbf n}
\def\bfo{\mathbf o}
\def\bfp{\mathbf p}
\def\bfq{\mathbf q}
\def\bfr{\mathbf r}
\def\bfs{\mathbf s}
\def\bft{\mathbf t}
\def\bfu{\mathbf u}
\def\bfv{\mathbf v}
\def\bfw{\mathbf w}
\def\bfx{\mathbf x}
\def\bfy{\mathbf y}
\def\bfz{\mathbf z}

\def\bfA{\mathbf A}
\def\bfB{\mathbf B}
\def\bfC{\mathbf C}
\def\bfD{\mathbf D}
\def\bfE{\mathbf E}
\def\bfF{\mathbf F}
\def\bfG{\mathbf G}
\def\bfH{\mathbf H}
\def\bfI{\mathbf I}
\def\bfJ{\mathbf J}
\def\bfK{\mathbf K}
\def\bfL{\mathbf L}
\def\bfM{\mathbf M}
\def\bfN{\mathbf N}
\def\bfO{\mathbf O}
\def\bfP{\mathbf P}
\def\bfQ{\mathbf Q}
\def\bfR{\mathbf R}
\def\bfS{\mathbf S}
\def\bfT{\mathbf T}
\def\bfU{\mathbf U}
\def\bfV{\mathbf V}
\def\bfW{\mathbf W}
\def\bfX{\mathbf X}
\def\bfY{\mathbf Y}
\def\bfZ{\mathbf Z}

\def\frD{\mathfrak D}
\def\frL{\mathfrak L}

\def\bfw{\mathbf w}

\def\R {{\mathbb R }}
 \def\C {{\mathbb C }}
 \def\Z{{\mathbb Z}}
 \def\H{{\mathbb H}}
\def\K{{\mathbb K}}
\def\N{{\mathbb N}}
\def\Q{{\mathbb Q}}
\def\A{{\mathbb A}}

\def\T{\mathbb T}
\def\P{\mathbb P}

%\def\G{\mathbb G}

\def\cD{\EuScript D}
\def\cL{\mathscr L}
\def\cK{\EuScript K}
\def\cM{\EuScript M}
\def\cN{\EuScript N}
\def\cR{\EuScript R}
\def\cW{\EuScript W}
\def\cY{\EuScript Y}
\def\cF{\EuScript F}

\def\bbA{\mathbb A}
\def\bbB{\mathbb B}
\def\bbD{\mathbb D}
\def\bbE{\mathbb E}
\def\bbF{\mathbb F}
\def\bbG{\mathbb G}
\def\bbI{\mathbb I}
\def\bbJ{\mathbb J}
\def\bbL{\mathbb L}
\def\bbM{\mathbb M}
\def\bbN{\mathbb N}
\def\bbO{\mathbb O}
\def\bbP{\mathbb P}
\def\bbQ{\mathbb Q}
\def\bbS{\mathbb S}
\def\bbT{\mathbb T}
\def\bbU{\mathbb U}
\def\bbV{\mathbb V}
\def\bbW{\mathbb W}
\def\bbX{\mathbb X}
\def\bbY{\mathbb Y}

\def\kappa{\varkappa}
\def\epsilon{\varepsilon}
\def\phi{\varphi}
\def\le{\leqslant}
\def\ge{\geqslant}

\def\B{\mathrm B}

\def\la{\langle}
\def\ra{\rangle}

\def\F{{}_2F_1}
\def\FF{{}_2^{\vphantom \C}F_1^\C}

\def\vp{^{\vphantom \C}}
\def\G{\vp G}

\newcommand{\dd}[1]{\,{\rm d}\,{\overline{\overline{#1}}} }

\def\lambdA{{\boldsymbol{\lambda}}}
\def\alphA{{\boldsymbol{\alpha}}}
\def\betA{{\boldsymbol{\beta}}}
\def\mU{{\boldsymbol{\mu}}}
\def\sigmA{{\boldsymbol{\sigma}}}

\def\1{\boldsymbol{1}}
\def\2{\boldsymbol{2}}

\def\SL{\mathrm{SL}}
\def\GL{\mathrm{GL}}
\def\even{\mathrm{even}}
\def\SO{\mathrm{SO}}

\def\B{\mathrm{B}}
% Root-approved finite author-wrapper index rendering. No generic fourIdx binding.
\makeatletter
\newcount\nativeIndexCount
\newcommand{\nativeIndexCountOne}{\ifmeasuring@\else\global\advance\nativeIndexCount by1\fi}
\makeatother
\ExplSyntaxOn
\cs_new_protected:Npn \native_index_error: {\PackageError{native-index}{Unobserved~index~wrapper~argument}{Only~the~fifteen~inventoried~author~calls~are~supported}}
\cs_new_protected:Npn \native_hyp_allowed:nn #1#2 {
 \str_if_eq:nnTF {#1}{r}{\str_if_eq:nnF{#2}{s}{\native_index_error:}}{
  \str_if_eq:nnTF{#1}{p}{\str_if_eq:nnF{#2}{q-1}{\str_if_eq:nnF{#2}{q}{\native_index_error:}}}{\native_index_error:}}
}
\cs_new_protected:Npn \native_complex_allowed:nn #1#2 {\str_if_eq:nnF{#1}{p}{\native_index_error:}\str_if_eq:nnF{#2}{q}{\native_index_error:}}
\cs_new_eq:NN \nativeHypAllowed \native_hyp_allowed:nn
\cs_new_eq:NN \nativeComplexAllowed \native_complex_allowed:nn
\cs_new_eq:NN \nativeIndexError \native_index_error:
\ExplSyntaxOff
\renewcommand{\hyp}[2]{\nativeHypAllowed{#1}{#2}\nativeIndexCountOne\mathord{\sideset{_{#1}}{_{\!#2}}{\mathop{F}}}}
\renewcommand{\Kc}[2]{\nativeComplexAllowed{#1}{#2}\nativeIndexCountOne\mathord{\sideset{_{#1}}{^{\C}_{\!#2}}{\mathop{\EuScript K}}}}
\renewcommand{\Gc}[2]{\nativeComplexAllowed{#1}{#2}\nativeIndexCountOne\mathord{\sideset{_{#1}}{^{\C}_{\!#2}}{\mathop{G}}}}
\renewcommand{\hypc}[2]{\nativeIndexError}

\AtEndDocument{\ifnum\nativeIndexCount=15\else\PackageError{native-index}{Expected exactly 15 inventoried calls}{Do not extend the finite resolver}\fi}

\makeatletter\newlabel{eq:beta}{{1.4}{}{Original source locator}{}{}}\makeatother
\begin{document}
\begin{center}\Large For q>p and noncoalescing positive-real-part paired parameters satisfying (1.13), the complex-field Barnes–Ismagilov integral equals the finite quadratic hypergeometric expression Sigma\_+\textasciicircum{}C\end{center}
\noindent\textbf{Stable ID:} 1575\quad\textbf{Author:} Yury A. Neretin\par
\noindent\textbf{Work:} Barnes--Ismagilov Integrals and Hypergeometric Functions of the Complex Field\par
\noindent\textbf{Source:} \url{https://sigma-journal.com/2020/072/}\par
\noindent\textbf{Version:} Official SIGMA 16 (2020) article 072, DOI 10.3842/SIGMA.2020.072; journal PDF and native TeX acquired 2026-09-30, exact bytes pinned by SHA256\par
\noindent\textbf{Rights:} CC BY 4.0. Authors retain copyright. \url{https://creativecommons.org/licenses/by/4.0/}. Reuse and typesetting adaptation; original language and mathematical content preserved.\par
\noindent\textbf{Difficulty (prior selection preserved):} {\CJKfont H2: The central obligation is to evaluate an integral with a simultaneous discrete Mellin index and continuous contour variable, whose paired Gamma poles must be matched across both indices. The proof derives their joint lattice parametrization, computes complex Gamma residues and reorganizes them into a finite sum of products of ordinary hypergeometric series with a convergence justification. This exceeds a routine single-variable residue exercise and does not merely apply the classical Meijer formula without the new discrete-index calculation.}\par
\medskip\noindent\textbf{Editorial boundary note (not author text):} Selected exactly the q>p,positive-real-part paired-parameter,noncoalescing coefficient-regular branch of author-main1.2(a), for z in C*. Full independent primed/unprimed pair and double-power conventions, ordinary hypergeometric/Pochhammer notation, measure/Mellin conventions, complex Gamma definitions/shifts/residues/asymptotics, every1.12–1.13 convergence hypothesis, kernel/imaginary-axis integral, exact Sigma-plus formula1.18 and complete new two-index residue-lattice/factorization/absolute convergence/local uniformity argument are retained. The precise classical Meijer Theorem18 is stated background; complete current Gamma asymptotics and pole separation proof remain editable. Other q<p/balanced branches remain literal statement context only; their analogous right-pole proof, further analytic continuation, differential systems and convolutions are unselected. Original1.8a prints (a-a), pole lemma900 prints Im sigma<0 whereas its statement870 uses Re, and the displayed residue expansion934 has its original Pochhammer-parenthesis form; all literal native formulas remain unchanged. No editor formula/sign/index correction or new proof is inserted. Counter resets preserve publisher numbering after unselected introductory examples and parameter-continuation paragraph; full native source stays privately archived. Root-approved finite wrapper-only AMS sideset typography for exactly15 inventoried hyp/Kc/Gc calls with exact argument whitelist and actual nonmeasuring count15; no generic fourIdx binding, no source call/argument/base-glyph edits, unused hypc and unknown arguments fail. Original fouridx import/macros and failed build remain hash-bound privately; unsafe official package retrieval was stopped without an alternative route.\par
\noindent\textbf{External original reference eq:beta:} Original complex beta integral source265–270, physical2, explicitly established in Gelfand/Graev/Vilenkin and Gelfand/Graev/Retakh. The unselected beta example is referenced only in the retained Gamma-function remark, not used as missing new core.\par
\medskip\hrule\medskip\noindent\textbf{Author text begins below.}\par
\setcounter{section}{1}\setcounter{subsection}{1}\setcounter{equation}{5}\setcounter{footnote}{1}
% Original source lines 321--547
\subsection{Notation}
Denote by $\Z_+$ the set of integers $\ge 0$, by $\Z_-$ the set of integers $\le 0$.

\subsubsection{Hypergeometric functions.}
Let $r\le s+1$.
Let $a_1,\dots,a_r\in\C$ and $b_1,\dots,b_s\in\C\setminus\Z_-$.
Generalized hypergeometric functions are defined by
\begin{equation}
\hyp{r}{s}\left[\genfrac{}{}{0pt}{}{a_1,\dots,a_r}{b_1,\dots,b_r};z\right]
:=\sum_{m=0}^\infty
\frac{(a_1)_m\cdots (a_r)_m}{m! (b_1)_m\cdots (b_s)_m}\cdot z^n,
\label{eq:rFs}
\end{equation}
where $(a)_m:=a(a+1)\cdots(a+m-1)$ is the Pochhammer symbol.
For $r\le s$ the radius of convergence is $\infty$, for $r=s+1$ it is $1$.
For $r>s+1$ the series diverges.

\subsubsection{Notation for lists}
Let $a_1, \dots, a_p\in \C$ and $h\in\C$.
We use the following notation for lists,
\begin{gather*}
(a):= a_1, \dots, a_p;\\
(a)+h:= a_1+h, \dots, a_p+h;\\
(a)_{\setminus j}:= a_1,\dots,a_{j-1},a_{j+1},\dots,a_p.
\end{gather*}
In particular, we denote \eqref{eq:rFs}
by $\hyp{r}{s}\big[\genfrac{}{}{0pt}{}{(a)}{(b)};z\big]$.

\subsubsection{Double powers}
Denote by $\C^\times$ (resp.\ $\R^\times_+$) the multiplicative group of $\C$
(resp.\ the multiplicative group of positive reals).
By $\T\subset \C^\times$ we denote the subgroup
$|z|=1$, we have $\T\simeq\R/2\pi \Z$,
\begin{equation*}
\C^\times= \T\times \R^\times_+,
%\label{eq:times1}
\end{equation*}
and $\R^\times_+$ is isomorphic to the additive group $\R$ of real numbers.

The dual group $(\C^\times)^*$, i.e., the group of all homomorphisms
$\C^\times\to \T$, is
\begin{equation*}
(\C^\times)^*= (\T)^*\times (\R^\times_+)^*\simeq \Z\times \R.
%\label{eq:times2}
\end{equation*}

Denote by $\Lambda_\C$ the set of pairs
$a\big| a'$ such that $a$, $a'\in\C$ and $a-a'\in\Z$.
We write such pairs by
\[
a\big|a'=\tfrac{k+\sigma}2\big|\tfrac{-k+\sigma}2,
\qquad\text{where $k\in\Z$, $\sigma\in\C$}.
\]
By $\Lambda\subset \Lambda_\C$ we denote the subset
consisting of $a|a'$ satisfying $a'+\ov a=0$,
\[
a\big|a'=
\tfrac{k+{\rm i}s}2\big|\tfrac{-k+{\rm i}s}2\in \Lambda,\qquad\text{where $s\in \R$}.
\]

Define the following functions on $\C^\times$:
\[
z^{a|a'}:= z^a \ov z^{\,a'}:= (z/\ov z)^{k/2} |z|^{\sigma}.
\]
Notice that such functions are precisely all homomorphisms
$\C^\times\to \C^\times$.
For $a|a'\in \Lambda$ we get homomorphisms $\C^\times\to \T$,
\[
\Big|z^{\tfrac{k+{\rm i}s}2\big|\tfrac{-k+{\rm i}s}2} \Big|=1.
\]

\subsubsection{The complex Mellin transform}
Denote by $\dd t$ the Lebesgue measure on $\C$,
\[
\dd t:= {\rm d}\Re t\, {\rm d}\Im t.
\]

We define the complex Mellin transform $\cM$ as the Fourier transform on the group $\C^\times$.
For a~function $f$ on $\C^\times$ we define the Mellin transform as a function
on $\Lambda_\C$ defined by
\[
g\big(a\big|a'\big)=
\cM f\big(a\big|a'\big):=\frac 1{2\pi}
\int_{\C^\times} t^{a|a'} f(t)\frac {\dd t}{t^{1|1}},
\]
the factor ${\dd t}/{t^{1|1}}$ is the $\C^\times$-invariant measure on $\C^\times$.

The Mellin transform is a unitary operator from
$L^2\big(\C^\times,{\dd t}/{t^{1|1}} \big)$
to $L^2$ on $\Lambda\simeq \Z\times\R$, the inversion formula is
\begin{equation*}
f(t)=\frac 1{2\pi}\sum_{k\in\Z} \int_\R g\big(\tfrac{k+{\rm i}s}2\big|\tfrac{-k+{\rm i}s}2\big) t^{{\tfrac{-k-{\rm i}s}2\big|\tfrac{k-{\rm i}s}2} }\,{\rm d}s.
%\label{eq:inversion}
\end{equation*}

The convolution on the group
$\C^\times$ is defined by formula
\[ f*g(z):=\int_\C f(t)g(z/t)\,\frac{\dd t}{t^{1|1}}.
\]
The Mellin transform sends convolutions to products,
\[
\cM (f*g)(z)=2\pi\, \cM f(z)\cdot \cM g(z).
\]

\subsection{The Gamma-function of the complex field}\label{ss:Gamma}
Following Gelfand, Graev and Retakh \cite{GGR},
we define the Gamma-function of the complex field as
a function on $\Lambda_\C$ by\footnote{The definition slightly
differs from a definition from \cite{GGR}, which was used in
\cite{MN,Ner-doug}. We write $\Im t$ instead of $\Re t$. For this reason
the factor $i^{a-a'}$ from \cite{GGR,MN} disappears.}
\begin{align}
\Gamma^\C\big(a\big|a'\big)&:=
\frac 1\pi \int_\C t^{a-1|a'-1} {\rm e}^{2{\rm i} \Im t}\dd ta \notag\\
&\hphantom{:}=
\frac{1}{\pi} \lim_{r\to\infty} \int_{|t|\le r} t^{a-1|a'-1}
{\rm e}^{2{\rm i} \Im t}\dd t
= \frac{\Gamma(a)}{\Gamma(1-a')}
=\frac{(-1)^{a-a'}\Gamma(a')}{\Gamma(1-a)} \notag \\
&\hphantom{:}= \frac{1}{\pi}\,
\Gamma(a)\Gamma(a')\sin\pi a'
= \frac{(-1)^{a-a'}}{\pi} \,\Gamma(a)\Gamma(a')\sin\pi a. \label{eq:Gamma}
\end{align}
The integral conditionally converges if $0<\Re(a+a')<1$
and diverges otherwise. Clearly, the right hand side is meromorphic
in the whole $\Lambda_\C$.

\begin{Remark} In particular, we can write \eqref{eq:beta} as
\[
\B^\C(\alpha|\alpha',\beta|\beta')=
\frac{\Gamma^\C(\alpha|\alpha')\,\Gamma^\C(\beta|\beta')}
{\Gamma^\C(\alpha+\beta|\alpha'+\beta')}.
\]
\end{Remark}

It is easy to see that
\begin{subequations}
\begin{gather}
\Gamma^\C\big(a|a'\big)=(-1)^{a-a} \Gamma^\C\big(a'|a\big),\label{eq:symmetry} \\
\Gamma^\C\big(a|a'\big)\, \Gamma^\C\big(1-a|1-a'\big) =(-1)^{a-a'},\label{eq:gamma-complement} \\
\Gamma^\C\big(a+m\big|a'+m'\big) =(-1)^{m'} \Gamma^\C\big(a\big|a'\big)(a)_m(a')_{m'}, \label{eq:shift1} \\
 \Gamma^\C\big(a-m\big|a'-m'\big) =\frac{(-1)^{m} \Gamma^\C\big(a\big|a'\big)}{(1-a)_m(1-a')_{m'}}, \label{eq:shift2} \\
\prod_{j=0}^{k-1}\Gamma^\C\big(a+\tfrac jk\big|a'+\tfrac jk \big)=\Gamma^\C\big(ka\big|ka'\big)k^{1-(a+a')k},\label{eq:gamma-gauss}
\end{gather}
\end{subequations}
where $m\in\N$, $k\in\N$.
Values of $\Gamma^\C$ at integer points are
\begin{gather*}
 \Gamma^\C(k|l) = 0,\qquad \text{for $k$, $l\in\N$;} \\
 \Gamma(m|-k) =\frac{(m-1)!(-1)^k}{k!},\qquad \text{for $m\in \N$, $k\in \Z_+$}.
\end{gather*}
For $m$, $m'\in\Z_+$ we have a pole at $(-m)|(-m')$, more precisely,
\begin{equation}
\res\limits_{\epsilon=0}
\Gamma^\C(-m+\epsilon|-m'+\epsilon)=\frac{(-1)^m}{m!\, m'!}.
\label{eq:res-gamma}
\end{equation}
Next (see Section~\ref{ss:asy} below),
for $a|a'\in\Lambda_\C$ and $\xi=\frac12(k+{\rm i}s)$
with $k\in\Z$, $s\in \R$ we have
\begin{equation}
\Gamma^\C\big(a+\xi|a'-\ov \xi\big)=
\exp\big\{2{\rm i}\Im (\xi\ln \xi-\xi)\big\}\cdot \xi^{a-\tfrac12\big|a'-\tfrac12} \big(1+O\big(|\xi|^{-1}\big)\big).
\label{eq:asy1}
\end{equation}
Therefore,
\begin{equation}
\big| \Gamma^\C\big(a+\xi|a'-\ov \xi\big)\big|\sim |\xi|^{\Re(a+a')-1}.
\label{eq:asy2}
\end{equation}

\subsection{Hypergeometric functions of the complex field}\label{ss:def}
Let $a_1|a'_1,\dots, a_q|a'_q$, $b_1|b'_1,\dots,b_p|b'_p\in \Lambda_\C$.
Temporarily, we assume that they satisfy the conditions:
\begin{equation}
\Re(a_\alpha+a_\alpha')>0,\qquad \Re(b_\beta+b_\beta')>0.
\label{eq:cond1}
\end{equation}
and
\begin{equation}
\sum_\alpha \Re(a_\alpha+a_\alpha')+\sum_\beta\Re(b_\beta+b_\beta')<p+q.
 \label{eq:cond2}
\end{equation}

We define the following function on $\Lambda_\C$:
\begin{gather}
\Kc{p}{q}\left[\genfrac{}{}{0pt}{}{(a|a')}{(b|b')};  \tfrac{k+\sigma}{2}\Bigl| \tfrac{-k+\sigma}{2}
\right]\nonumber\\
\qquad{} :=
 \prod_{\alpha=1}^q \Gamma^\C\Bigl(a_\alpha+\tfrac{k+\sigma}2\Bigl|a_\alpha'+\tfrac{-k+\sigma}2\Bigr)
 \prod_{\beta=1}^p \Gamma^\C\Bigl(b_\beta+\tfrac{-k-\sigma}2\Bigl|b_\beta'+\tfrac{k-\sigma}2\Bigr).
\label{eq:cK}
\end{gather}
Next, we define the hypergeometric functions of the complex field
as the contour integral:
\begin{equation}
\Gc{p}{q}\left[\genfrac{}{}{0pt}{}{(a|a')}{(b|b')};z\right]
:=\frac{1}{2\pi {\rm i}}
\sum_{k\in\Z} \int_{{\rm i}\R}
\Kc{p}{q}\left[\genfrac{}{}{0pt}{}{(a|a')}{(b|b')};\tfrac{k+\sigma}{2}\Bigl| \tfrac{-k+\sigma}{2}\right]
 z^{\frac{-k-\sigma}2\big| \frac{k-\sigma}2} \,{\rm d}\sigma.
 \label{eq:def}
\end{equation}
The integration is taken along the imaginary axis ${\rm i}\R$.
The condition \eqref{eq:cond2} provides the conditional convergence of
the integral in $\sigma$ and the absolute convergence of the series
(see Sections~\ref{ss:meijer} and \ref{ss:evaluation} below).
Under the stronger condition
\begin{equation}
\sum_\alpha \Re(a_\alpha+a_\alpha')+\sum_\beta\Re(b_\beta+b_\beta')<p+q-1
 \label{eq:cond3}
\end{equation}
all integrals in $\sigma$ convergence absolutely
(this follows from \eqref{eq:asy2}).

The expression \eqref{eq:cK} has poles (they are analyzed in
Section~\ref{ss:evaluation}) in $\sigma$ originating from two groups
of factors.
We call poles of the factors
$\Gamma^\C\big(a+\frac{k+\sigma}2\big|a'+\frac{-k+\sigma}2\big)$
\emph{left poles}, the poles of the factors
$\Gamma^\C\big(b+\frac{-k-\sigma}2\big|b'+\frac{k-\sigma}2\big)$
\emph{right poles}.
Under the conditions \eqref{eq:cond1} left poles are contained
in the left half-plane $\Re\sigma<0$,
right poles are contained in the right half plane $\Re\sigma>0$, and
the axis ${\rm i}\R$ separates these groups of poles.

\setcounter{equation}{17}
% Original source lines 571--617
\subsection{The statements of the paper}
For the same $(a|a')$, $(b|b')$ we define the expression
$\Sigma^\C_+(z)$ by
\begin{gather}
\Sigma^\C_+
\left[\genfrac{}{}{0pt}{}{(a|a')}{(b|b')};z\right]
:=2\sum_{j=1}^q z^{a_j|a_j'}\cdot \prod_{\beta}
\Gamma^\C\big(b_\beta+a_j\big|b_\beta'+a_j'\big)
\cdot \prod_{\alpha\ne j} \Gamma^\C\big(a_\alpha-a_j\big|a_\alpha'-a_j'\big) \nonumber\\
\qquad{}
\times
\hyp{p}{q-1}\left[\genfrac{}{}{0pt}{}{(b_\alpha+a_j)}
{(1-a_\alpha+a_j)_{\setminus j}};(-1)^q\,z\right]
\hyp{p}{q-1}\left[\genfrac{}{}{0pt}{}{(b_\alpha'+a_j')}
{(1-a_\alpha'+a_j')_{\setminus j}};(-1)^p\,\ov z\right].
 \label{eq:Xi}
\end{gather}
We also define the expression $\Sigma^\C_-(z)$ by
\[
\Sigma^\C_-
\left[\genfrac{}{}{0pt}{}{(a|a')}{(b|b')};z\right]
:=\Sigma^\C_+ \left[\genfrac{}{}{0pt}{}{(b|b')}{(a|a')};z^{-1}\right].
\]

\begin{Theorem}\label{th:1}
Let $(a|a')$, $(b|b')$ satisfy the conditions \eqref{eq:cond1} and~\eqref{eq:cond2}.
\begin{enumerate}\itemsep=0pt
\item[$(a)$] For $q>p$ we have
\[
\Gc{p}{q}\left[\genfrac{}{}{0pt}{}{(a|a')}{(b|b')};z\right]
=\Sigma^\C_+ \left[\genfrac{}{}{0pt}{}{(a|a')}{(b|b')};z\right]
\]
and for $q<p$ we have $_p\G^\C_q(z)=\Sigma^\C_-(z)$.
\item[$(b)$] For $p=q$
\[
\Gc{p}{q}(z)=
\begin{cases}
\Sigma^\C_+(z), & \text{if $|z|<1$}, \\
\Sigma^\C_-(z), & \text{if $|z|>1$}.
\end{cases}
\]
\end{enumerate}
\end{Theorem}

\begin{Remark} This statement and the main argument for its proof
are potentially contained in the paper by Ismagilov \cite[Lemma~2]{Ism2}.
\end{Remark}


% Original source lines 735--987
\section{Proofs}\label{section2}

\subsection[Asymptotics of $\Gamma^\C$]{Asymptotics of $\boldsymbol{\Gamma^\C}$}\label{ss:asy}
Let us derive formula \eqref{eq:asy1} for the asymptotics of the function $\Gamma^C(a+\xi|a'-\ov \xi)$.
First, let us verify that the right-hand side is single-valued
on $\Lambda$.
We must verify that the expression
\[
\exp\big\{ \xi\ln\xi - \ov\xi\ln\ov\xi\big\}, \qquad \xi=\tfrac{k+{\rm i}s}2
\]
is single valued.
Represent $\xi=r {\rm e}^{{\rm i}\phi}$. We transform our expression as
\begin{gather*}
\exp\big\{\xi(\ln(r)+{\rm i}\phi+2\pi {\rm i} n)-
\ov \xi(\ln(r)-{\rm i}\phi-2\pi {\rm i} n) \big\} \\
\qquad{}=\exp\big\{(\xi-\ov\xi) \ln r +(\xi+\ov\xi)({\rm i}\phi+2\pi {\rm i} n) \big\}.
\end{gather*}
We have $\xi+\ov\xi\in \Z$ and therefore the exponential is single valued.

Next, we apply the Stirling formula in the form
(see \cite[equation~(1.18(3))]{HTF1})
\[
\ln
\Gamma(c+z)=\big(z+c-\tfrac12\big)\ln z -z+\tfrac 12\ln(2\pi)+O\big(z^{-1}\big),
\qquad \text{where $|\arg (z)|<\pi-\epsilon$},
\]
and get
\begin{gather*}
\Gamma(a+\xi) \simeq \sqrt{2\pi}\exp
\big\{ \big(\xi+a-\tfrac12\big)\ln\xi-\xi \big\},\\
\Gamma(1-a'-\ov \xi) \simeq \sqrt{2\pi}\exp
\big\{ \big(\ov\xi-a'+\tfrac12\big)\ln\ov\xi-\ov\xi\big\}.
\end{gather*}
The ratio gives us \eqref{eq:asy1}.
An evaluation of the asymptotics in the sector $|\arg(-z)|<\pi-\epsilon$
gives the same result.

\begin{Corollary}\label{cor:L2}
Let $\Re(a_\alpha+a_\alpha')>0$, $\Re(b_\beta+b_\beta')>0$.
Denote
\[
\upsilon:=\sum (a_\alpha+a_\alpha')+\sum(b_\beta+b_\beta').
\]
\begin{enumerate}\itemsep=0pt
\item[$(a)$] If $\upsilon<p+q$, then $F(k,{\rm i}s):=
\Kc{p}{q}\big[\genfrac{}{}{0pt}{}{(a|a')}{(b|b')};\tfrac{k+{\rm i}s}{2}\bigl|;\tfrac{-k+{\rm i}s}{2}\big]$
tends to $0$ as $|k+{\rm i}s|$ tends to $\infty$.
\item[$(b)$]
If $\upsilon<p+q-1$, then for each $k$ the function $F(k,{\rm i}s)$ as a
function in $s$ is integrable.
Under the same condition $F(k,{\rm i}s)$ is contained in $L^2(\Lambda)$.
\item[$(c)$]
If $\upsilon<p+q-2$, then $F(k,{\rm i}s)$ is contained in $L^1(\Lambda)$.
\end{enumerate}
\end{Corollary}

\begin{Remark}
Formula \eqref{eq:asy1} also gives us asymptotics of
$\Kc{p}{q}$ on vertical lines $\sigma=h+{\rm i}s$. Indeed,
\[
\Gamma\big(a+\tfrac{k+h+{\rm i}s}{2}\big|a'+\tfrac{-k+h+{\rm i}s}{2}\big)=
\Gamma\big(a+\tfrac h2+\tfrac{k+{\rm i}s}{2}\big|a'+\tfrac h2+\tfrac{-k+{\rm i}s}{2}\big),
\]
and we can control integrability under shifts of the integration contour.
\end{Remark}

\subsection{Decomposition of Mellin--Barnes integrals in residues}\label{ss:meijer}
Now we start to prove Theorem~\ref{th:1}, i.e., to derive
the quadratic expressions of the func\-tions~$\Gc{p}{q}$ in terms of
the usual hypergeometric functions $\hyp{p}{q}$.

First, we write the definition \eqref{eq:def} of $\Gc{p}{q}$ in the form
\[
\sum_{k\in \Z} z^{-k/2\big|k/2} I_k(z), \qquad
I_k(z)=\frac{1}{2\pi {\rm i}}\int (\dots)|z|^{-\sigma}\,{\rm d}\sigma.
\]

The integrals $I_k(z)$ are special cases of
Mellin--Barnes integrals, i.e., integrals of the type
\begin{equation}
J^{A,B}_{C,D}(a,b,c,d;u)=\frac 1{2\pi {\rm i}}\int_L \frac{\prod\limits_{\alpha=1}^A\Gamma(a_\alpha+\sigma)
\prod\limits_{\beta=1}^B\Gamma(b_\beta-\sigma)}
{\prod\limits_{\gamma=1}^C\Gamma(c_\gamma+\sigma)\prod\limits_{\delta=1}^D\Gamma(d_\delta-\sigma)} u^{-s}
\, {\rm d}s,
\label{eq:J}
\end{equation}
where $L$ is a contour separating poles of factors
$\Gamma(a_\alpha+\sigma)$ and poles of $\Gamma(b_\beta-\sigma)$.
The behavior of such Mellin--Barnes integrals
(Meijer $G$-functions) was investigated by Meijer,
his results are exposed in \cite{BW, Luke,Mar}.
Under certain conditions $J(a,b,c,d;u)$ can be expressed in terms
of functions $\Sigma_\pm(z)$, where
$\Sigma_+(z)$ is the sum of residues of the integrand at poles of
factors $\Gamma(a_\alpha+\sigma)$, and $\Sigma_-(z)$
the sum of residues at poles of $\Gamma(b_\beta-\sigma)$.

In our case $A=D=q$, $B=C=p$, and $u=|z|$ is a positive real.
The contour $L$ coincides with the imaginary axis at infinity.
The asymptotics of the absolute value of the integrand on the imaginary axis
is
\[
\sim |\sigma|^{\sum\limits_{\alpha=1}^q\Re (a_\alpha-d_\alpha)+\sum\limits_{\beta=1}^p\Re(b_\beta- c_\beta)},
\]
see \cite[equation~(1.18(4))]{HTF1}.
One of the statements of \cite[Theorem~18]{Mar} implies that
\emph{if the integrand tends to zero on the imaginary axis,
then the integral $J(a,b,c,d;u)$ conditionally converges if
the integrand tends to $0$ at infinity for $u>0$ and is given by}
\begin{itemize}\itemsep=0pt
\item[--] \textit{$\Sigma_+(u)$ for $q>p$;}
\item[--] \textit{$\Sigma_-(u)$ for $q<p$;}
\item[--] \textit{for $p=q$ we have $\Sigma_+(u)$ if $0<u<1$ and $\Sigma_-(u)$ for $u>1$.}
\end{itemize}
If
$\sum\Re(a_\alpha-d_\alpha)+\sum\Re(b_\beta-c_\beta)<-1$,
then the absolute convergence is obvious.

In our case $a_\alpha$ is replaced by $a_\alpha+k/2$, $b_\beta$ by
$b_\beta-k/2$, and $c_\beta\mapsto 1-b'_\beta-k/2$,
$d_\alpha\mapsto 1-a'_\alpha+k/2$.
Therefore under our conditions \eqref{eq:cond1}, \eqref{eq:cond2}
all integrals $I_k$ converge.

\subsection{Evaluation of sums of residues}\label{ss:evaluation}

\begin{Lemma}
Fix $a|a'$ and consider the following family of functions:
\[
\gamma_k(\sigma):=\Gamma^\C\big(a+\tfrac {k+\sigma}2\big|a'+\tfrac
{-k+\sigma}2\big).
\]
Denote by $\Omega$ the sets of all points $(k,\sigma)\in \Z\times \C$
such that $\sigma$ is a pole of $\gamma_k(\sigma)$.
\begin{enumerate}\itemsep=0pt
\item[$(a)$] The set $\Omega$ is contained in the half-plane $\Re\sigma<0$ if and only
if $\Re(a+a')>0$.
\item[$(b)$] The set $\Omega$ consists of points
\begin{gather*}
k=-m+m'-a+a', \\
\sigma= -m-m'-a-a',
\end{gather*}
where $m$, $m'$ range in $\Z_+$.
\end{enumerate}
\end{Lemma}

\begin{proof} (a)
Let all poles be contained in the domain $\Re\sigma<0$.
Taking $k=a'-a$ we get
\[
\gamma_{a'-a}(\sigma)=\Gamma^\C
\big(a+\tfrac {a'-a+\sigma}2\big|a'+\tfrac {-a'+a+\sigma}2\big)=
\Gamma^\C\big(\tfrac{a+a'}2 +\tfrac{\sigma}{2}\big|\tfrac{a+a'}2
+\tfrac{\sigma}{2}\big).
\]
The point $\sigma= -(a+a')$ is a pole of $\gamma_{a'-a}(\sigma)$.
Therefore, $\Re(a+a')>0$.

Conversely, let $\Re(a+a')>0$.
Consider $(k,\sigma)\in \Omega$.
Then (see Section~\ref{ss:Gamma})
\begin{equation}
a+\tfrac{k+\sigma}2=-m\in \Z_-, \qquad a'+\tfrac{-k+\sigma}2=-m'\in \Z_-.
\label{eq:Omega}
\end{equation}
Therefore $a+a'+\sigma\in \Z_-$ and $\Im\sigma<0$.

(b) We solve the system \eqref{eq:Omega}.
\end{proof}

\begin{proof}[Evaluation of the sum of residues (the proof of Theorem~\ref{th:1})]
For definiteness assume that $p\le q$.
Let us write the residue $R(j;m,m')$ of the integrand \eqref{eq:def}
at a point
\[
k= -m+m'-a_j+a'_j,\qquad \sigma=-m-m'-a_j-a'_j.
\]
We have
\[
\tfrac{k+\sigma}{2}=-a_j-m,\qquad \tfrac{-k+\sigma}{2}=-a_j'-m',
\]
and
keeping in mind \eqref{eq:res-gamma} we get
\begin{gather*}
R(j;m,m')=\frac{2(-1)^m}{m!\,m'!}
 \prod_{\alpha\ne j}
\Gamma^\C\big(a_\alpha-a_j-m\big|a_\alpha'-a_j'-m'\big)\\
\phantom{R(j;m,m')=}{} \times
\prod_{\beta} \Gamma^\C\big(b_\beta+a_j+m\big|b_\beta'+a_j'+m'\big)
\times z^{a_j+m|a_j'+m'}.
\end{gather*}
Applying \eqref{eq:shift1} and \eqref{eq:shift2} we come to
\begin{gather*}
R(j;m,m')=2 \frac{(-1)^m}{m!\,m'!}
\prod_{\alpha\ne j} \frac{\Gamma^\C\big(
a_\alpha-a_j\big|a_\alpha'-a_j'\big)(-1)^{m}}
{(1-a_\alpha+a_j)_m(1-a_\alpha'+a_j')_{m'}} \\
\hphantom{R(j;m,m')=}{} \times
\prod_{\beta} \Gamma^\C\big(b_\beta+a_j\big|b_\beta'+a_j'\big)
b_\beta+a_j\big)_m(b_\beta'+a_j')_{m'} (-1)^{m'}
\times z^{a_j+m|a_j'+m'}.
\end{gather*}
Reordering factors, we get
\begin{gather*}
2 z^{a_j|a'_j}
\prod_{\alpha\ne j} \Gamma^\C\big(a_\alpha-a_j\big|a_\alpha'-a_j'\big)
\prod_{\beta}
\Gamma^\C\big(b_\beta+a_j\big|b_\beta'+a_j'\big) \\
\qquad{}\times
(-1)^{mq} \frac{z^m \prod\limits_{\beta}(b_\beta+a_j\big)_{m}}
{m!\prod\limits_{\alpha\ne j} (1-a_\alpha+a_j)_{m}}
 \times
(-1)^{m'p} \frac{\ov z^{m'} \prod\limits_{\beta}(b_\beta'+a_j'\big)_{m'}}
{m'!\prod\limits_{\alpha\ne j} (1-a_\alpha'+a_j')_{m'}}.
\end{gather*}

Formal calculation with series gives
\begin{gather*}
\sum_{m,m'} R(j,m,m')=
2 z^{a_j|a_j'}
 \prod_{\alpha\ne j} \Gamma^\C\big(a_\alpha-a_j\big|a_\alpha'-a_j'\big)
 \prod_{\beta}
 \Gamma^\C\big(b_\beta+a_j\big|b_\beta'+a_j'\big)
\\
\hphantom{\sum_{m,m'} R(j,m,m')=}{} \times
 \sum_m \frac{ \big((-1)^{q} z\big)^m \prod\limits_{\beta}(b_\beta+a_j\big)_m} {m!\prod\limits_{\alpha\ne j} (1-a_\alpha+a_j)_m}
 \times
 \sum_{m'} \frac{ \big((-1)^{p} \ov z\big)^{m'} \prod\limits_{\beta}(b_\beta+a_j\big)_{m'}}
 {m'!\prod\limits_{\alpha\ne j} (1-a_\alpha'+a_j')_{m'}} .
 \end{gather*}
We get a product of two hypergeometric series whose radius of convergence
is $\infty$ if $q>p$ and $1$ if $q=p$.
They are absolutely convergent in the disk of convergence,
therefore the last identity really takes place.
Therefore
\[
\sum_j \sum_{m,m'} R(j,m,m')=\Sigma^\C_+ (z).
\]
On the other hand the absolute convergence allows us to write
 \begin{equation}
 \sum_{m,m'} R(j,m,m')=\sum_{k\in\Z}\,\,
\sum_{m,m'\colon m-m'=k} R(j,m,m')=\sum_{k\in\Z} (z/\ov z)^{k/2} I_k.
 \label{eq:sum-R}
 \end{equation}
Thus we proved the coincidence of $\Gc{p}{q}$ and $\Sigma_+^\C$.
The summation of residues at the right poles of
$_p^{\vphantom{\C}}\cK^\C_q$ is similar.

It is important that the convergence of the series $\sum\limits_{k\in\Z}$
in~\eqref{eq:sum-R} is locally uniform in the parameters
$a_\alpha$, $b_\beta$ near any point $(a_0)\in\C^q$, $(b_0)\in\C^p$,
for which the coefficients at $z^l\ov z^{l'}$ are well-defined.
\end{proof}
\begin{thebibliography}{99}
\bibitem[3]{BW}
Beals R., Wong R., Special functions and orthogonal polynomials,
 \textit{Cambridge Studies in Advanced Mathematics}, Vol.~153, \href{https://doi.org/10.1017/CBO9781316227381}{Cambridge
 University Press}, Cambridge, 2016.

\bibitem[13]{HTF1}
Erd\'elyi A., Magnus W., Oberhettinger F., Tricomi F.G., Higher transcendental
 functions, {V}ol.~{I}, McGraw-Hill Book Company, Inc., New York~-- Toronto~--
 London, 1953.

\bibitem[16]{GGR}
Gel'fand I.M., Graev M.I., Retakh V.S., Hypergeometric functions over an
 arbitrary field, \href{https://doi.org/10.1070/RM2004v059n05ABEH000771}{\textit{Russian Math. Surveys}} \textbf{59} (2004), 831--905.

\bibitem[24]{Ism2}
Ismagilov R.S., Racah operators for principal series of representations of the
 group ${\rm SL}(2,{\mathbb C})$, \href{https://doi.org/10.1070/SM2007v198n03ABEH003840}{\textit{Sb. Math.}} \textbf{198} (2007),
 369--381.

\bibitem[30]{Luke}
Luke Y.L., The special functions and their approximations, {V}ol.~{I},
 \textit{Mathematics in Science and Engineering}, Vol.~53, Academic Press, New
 York~-- London, 1969.

\bibitem[31]{Mar}
Marichev O.I., Handbook of integral transforms of higher transcendental
 functions. Theory and algorithmic tables, \textit{Ellis Horwood Series: Mathematics
 and its Applications}, \href{https://doi.org/10.1002/zamm.19840640608}{Ellis Horwood Ltd.}, Chichester,
 1983 (translated from: A method of calculating integrals from special functions
(theory and tables of formulas), Nauka i Tekhnika, Minsk, 1978).

\bibitem[33]{MN}
Molchanov V.F., Neretin Yu.A., A pair of commuting hypergeometric operators on
 the complex plane and bispectrality, \textit{J.~Spectr. Theory}, {t}o appear,
 \href{https://arxiv.org/abs/1812.06766}{arXiv:1812.06766}.

\bibitem[39]{Ner-doug}
Neretin Yu.A., An analog of the {D}ougall formula and of the de
 {B}ranges--{W}ilson integral, \href{https://doi.org/10.1007/s11139-019-00218-0}{\textit{Ramanujan~J.}}, {t}o appear,
 \href{https://arxiv.org/abs/1812.07341}{arXiv:1812.07341}.

\end{thebibliography}
\end{document}
